\documentclass[10pt,a4paper]{amsart}
\usepackage[margin=19mm]{geometry}
\usepackage{amsmath,amssymb,mathtools}
\usepackage[scr=rsfs,cal=euler]{mathalfa}
\usepackage{xfrac}
\usepackage[unicode,hidelinks,bookmarks=false]{hyperref}
\newcommand{\ie}{i.e.\ }
\newcommand{\cf}{cf.\ }
\newcommand{\resp}{respectively\ }
\newcommand{\Subsection}{\subsection}
\providecommand{\nobreakdash}{\nobreak\hbox{-}}
\setlength{\emergencystretch}{2em}
\multlinegap0pt
\usepackage{amssymb}
\usepackage{mathtools}
\usepackage[shortlabels]{enumitem}
\newtheorem{lemma}{Lemma}
\newcommand{\Prob}{\mathbb P}
\newcommand{\Pts}{\operatorname{Pts}}
\newcommand{\Unif}{\operatorname{Unif}}
\newcommand{\qbinom}[2]{\genfrac{[}{]}{0pt}{}{#1}{#2}_q}
\title{Default-Distance Entropy and Metric Dimension in Finite Geometries}
\author{Maximiliano Vazquez}
\date{Source: arXiv:2608.27788v1 (CC BY 4.0)}
\begin{document}
\maketitle

\noindent\emph{Source and licence.} Default-Distance Entropy and Metric Dimension in Finite Geometries. Version arXiv:2608.27788v1 (CC BY 4.0), distributed by arXiv under the Creative Commons Attribution 4.0 International licence, http://creativecommons.org/licenses/by/4.0/, as recorded in the arXiv licence metadata for this exact version and shown on the abstract page https://arxiv.org/abs/2608.27788v1. Full licence notice: \texttt{licenses/arxiv-2608.27788-CC-BY-4.0.txt}.\par\smallskip

\noindent\textit{Editorial scope.} Counted unit: the article's lemma lem:dual-both (source0.tex lines 1296--1304). The article's complete original proof of it is delivered with it (source0.tex lines 1306--1357, one \texttt{proof} environment of 52 lines). No mathematics is written, repaired or re-derived: the statement and the proof are the article's own lines, in the article's own notation, and no retained line is substituted or re-flowed.  Retained uncounted as the source-local context the proof needs: lines 1167--1170.  Nothing was omitted from any retained span, so the package carries no omission record.  No retained line contains a citation, so the wrapper carries no bibliography and no bibliography entry.  No retained line contains a figure environment, an image-inclusion command or drawing code, so no image is delivered and the package declares no asset dependency.  Editorial formatting only, in addition: an AMS \texttt{amsart} preamble that restates the article's own declarations and macros used by the retained lines, namely the environments \texttt{lemma} on separate counters, together with the standard packages those lines need.  The retained environments print in this document as Lemma 1 in order of appearance, whereas the article prints the same items with the numbering of its own full document.  The complete native source of the article as distributed in the arXiv e-print for this exact version is bundled unchanged as \texttt{sources/208741/source0.tex}.
\begin{equation}\label{eq:dual-residue-count}
        N_t=\prod_{j=1}^t(q^{e+j-1}+1),
        \qquad N_0=1.
\end{equation}

\begin{lemma}\label{lem:dual-both}
Assume $d\geq2$ and $e>0$ are fixed.  If $U,V\in\mathcal G_d$ are distinct
and $Z\sim\Unif(\mathcal G_d)$, then, uniformly over all $U\neq V$,
\[
        \Prob(Z\cap U\neq0,\ Z\cap V\neq0)
        =O_{d,e}(q^{-e-1}+q^{-2e})
        =o(q^{-e}).
\]
\end{lemma}

\begin{proof}
Suppose that $Z$ meets both $U$ and $V$.  Choose projective points
$p\in\Pts(Z\cap U)$ and $p'\in\Pts(Z\cap V)$.  If $p=p'$, then $Z$ 
contains
a point of $U\cap V$.  If $p\neq p'$, then $p$ and $p'$ lie in the 
singular
generator $Z$, so they span a singular line.  We bound these two cases
separately.

Since $U\neq V$, the vector dimension of $U\cap V$ is at most $d-1$.
Consequently,
\[
        |\Pts(U\cap V)|\leq\qbinom{d-1}{1}=O_d(q^{d-2}).
\]
A fixed singular point lies in $N_{d-1}$ generators.  By
\eqref{eq:dual-residue-count}, the probability that $Z$ contains at least 
one
common point of $U$ and $V$ is therefore at most
\[
        O_d(q^{d-2})\frac{N_{d-1}}{N_d}
        =O_d(q^{d-2})\frac{1}{q^{e+d-1}+1}
        =O_{d,e}(q^{-e-1}).
\]

It remains to count ordered pairs $(p,p')\in\Pts(U)\times\Pts(V)$ that 
span
a singular line.  If $p\in U\cap V$, then there are $O_d(q^{d-2})$ choices
for $p$.  We may overcount by allowing all $O_d(q^{d-1})$ choices for
$p'\in\Pts(V)$, contributing
$O_d(q^{2d-3})$ pairs.  Now suppose $p\in U\setminus V$.  The singular 
points
of $V$ collinear with $p$ are precisely the points of $V\cap p^\perp$.
This is a proper subspace of $V$: if $V\subseteq p^\perp$, then
$\langle V,p\rangle$ would be a singular subspace of vector dimension
$d+1$, contradicting the maximality of $V$.  Hence
\[
        |\Pts(V\cap p^\perp)|=O_d(q^{d-2}).
\]
There are $O_d(q^{d-1})$ choices for $p$, so the number of ordered 
collinear
pairs is again $O_d(q^{2d-3})$, uniformly in $U$ and $V$.

Each singular line lies in $N_{d-2}$ generators.  A union bound over the
ordered pairs therefore gives
\[
        O_d(q^{2d-3})\frac{N_{d-2}}{N_d}
        =O_d(q^{2d-3})
          \frac{1}{(q^{e+d-2}+1)(q^{e+d-1}+1)}
        =O_{d,e}(q^{-2e}).
\]
Adding the common-point and singular-line contributions proves the lemma.
\end{proof}

\end{document}
